\begin{decisione}{D0002}{Come configurare la coda}{Proposta}{23/10/2026}

\contesto
Con la coda, un ordine a cui abbiamo risposto ``ricevuto'' esiste solo dentro RabbitMQ finché un worker non lo salva. Se la coda o un worker si spengono in quel momento, l'ordine non deve sparire.

\scelta
\begin{itemize}
  \item \textbf{Una sola istanza} di RabbitMQ, fissata a worker-1, con i dati su disco: è volutamente un punto unico di guasto (F3).
  \item \textbf{Coda e messaggi durevoli}: sopravvivono al riavvio di RabbitMQ.
  \item \textbf{Conferma all'API} (\textit{publisher confirm}): l'API risponde 202 solo quando RabbitMQ conferma di aver scritto il messaggio.
  \item \textbf{Conferma del worker} (\textit{ack}) solo \emph{dopo} il salvataggio nel database: se il worker si spegne prima, RabbitMQ consegna l'ordine a un altro worker.
  \item \textbf{Un ordine alla volta per worker} (\textit{prefetch} = 1): il lavoro si divide in modo equo, e un worker che si spegne lascia in sospeso al massimo un ordine.
  \item \textbf{Coda dei falliti} (DLQ): raccoglie gli ordini che non si riesce a salvare dopo le riprove (D0003). A fine esperimento vengono rielaborati e contati a parte.
\end{itemize}

\conseguenze
Un ordine confermato al cliente non si perde per lo spegnimento della coda (F3) o dei worker (F2). Il prezzo è che lo stesso ordine può arrivare due volte (dopo un worker spento tra salvataggio e conferma): per questo serve D0004. Scrivere su disco e aspettare la conferma rende l'API un po' più lenta, ed è parte di ciò che misuriamo in RQ2.

\end{decisione}
